\section{Overview}
\label{sec:overview}

Before the formal development we walk through the four running examples,
written in Rocq-like syntax on the left and in the guarded translation on the
right.  Throughout, $\T(t,u)$ is the typing predicate of the target logic
(``$t$ has type $u$''), and predicates and functions are translated to
first-order symbols applied to their type and term arguments.

\paragraph{A covered guard.}
The premise $\forall x{:}\mathit{nat}.\ \mathit{Even}\,x \imp
\mathit{Even}\,(S\,(S\,x))$ translates to
\begin{equation}
  \label{eq:even-guarded}
  \forall x.\ \T(x,\mathit{nat}) \imp \mathit{Even}(x) \imp
  \mathit{Even}(S(S(x))).
\end{equation}
The guard is redundant.  In the intended interpretation, an atom
$\mathit{Even}(t)$ is true only if the source application
$\mathit{Even}\;t$ is a well-formed \emph{and} provable proposition; since
$\mathit{Even} : \mathit{nat} \arr \PProp$, well-formedness forces
$t : \mathit{nat}$.  We say the atom $\mathit{Even}(x)$ \emph{covers} the
guard $\T(x,\mathit{nat})$: under any valuation sending $x$ outside
$\mathit{nat}$, the atom is false and the body of \eqref{eq:even-guarded} is
vacuously true.  Hence \eqref{eq:even-guarded} and
$\forall x.\ \mathit{Even}(x) \imp \mathit{Even}(S(S(x)))$ have the same
truth value in the intended model --- an \emph{equivalence}, not merely an
implication, so the rewrite is legitimate at any polarity.  Moreover the
equivalence is provable in pure first-order logic from the
\emph{typehood-completion axiom}
\[
  \forall x.\ \mathit{Even}(x) \imp \T(x,\mathit{nat}),
\]
which is itself true in the intended model.  This observation drives the
architecture of \cref{sec:guards}: guard omission is a logical equivalence
relative to a stock of completion axioms, and the semantic content of the
soundness proof is concentrated in the single fact that completion axioms
hold in the intended model.

\paragraph{An inversion axiom.}
For the inductive predicate $\mathit{le}$ the translation emits an inversion
axiom
\[
  \forall n\, m.\ \T(n,\mathit{nat}) \imp \T(m,\mathit{nat}) \imp
  \mathit{le}(n,m) \imp
  \bigl( m = n \vee \exists k.\ \T(k,\mathit{nat}) \wedge \mathit{le}(n,k)
  \wedge m = S(k) \bigr).
\]
Both universal guards are covered by the hypothesis atom
$\mathit{le}(n,m)$ and may be dropped.  Note that the criterion must look
\emph{past} the covering atom, not into the disjunction: the disjunct
$m = n$ contains $m$ ``naked'' in an equation, and equations force no typing
(junk equals junk).  The recursive judgment $\jt$ of \cref{sec:guards} is
therefore satisfied at the implication whose antecedent is
$\mathit{le}(n,m)$, regardless of the shape of what follows.  The
\emph{existential} conjunct $\T(k,\mathit{nat})$ is different: it is a fact
provided to the prover, not an obligation; it may be dropped either as a
weakening at a positive position of a premise
(\cref{def:admissible,thm:guards-sound}) or, here, as an equivalence via the
dual junk-\emph{falsity} judgment, because the remaining conjunct
$\mathit{le}(n,k)$ covers $k$ (\cref{thm:guard-equivalence}); the tightness
example of \cref{sec:tightness} shows that licensing existential guard
omission by the universal (junk-truth) rule instead would be unsound.

\paragraph{A guard that must stay, and one covered by it.}
The lemma $\mathit{app\_nil\_r} : \forall A{:}\mathsf{Type}.\ \forall
l{:}\mathit{list}\,A.\ l \mathbin{+\!\!+} \mathit{nil} = l$ translates to
\[
  \forall \alpha.\ \T(\alpha,\TypeC) \imp \forall l.\
  \T(l, \mathit{list}(\alpha)) \imp
  \mathit{app}(\alpha, l, \mathit{nil}(\alpha)) = l .
\]
The guard on $l$ is \emph{not} omissible: the body is a bare equation, no
atom covers $l$, and dropping the guard yields a sentence that is false in
the intended model (the equation fails for junk $l$, where the left-hand
side is a stuck application).  The guard on $\alpha$, however, \emph{is}
omissible: the retained guard $\T(l,\mathit{list}(\alpha))$ itself covers
it, because typing atoms whose type side is a rigid application of an
inductive type former force well-formedness of the former's arguments
($\T(l,\mathit{list}(\alpha))$ can hold only when $\alpha$ denotes a type).
Covers may thus come from ordinary hypothesis atoms and from \emph{retained}
guards; since a guard that serves as a cover must itself not be dropped, the
omissible set is computed as a greatest fixpoint, and
\cref{prop:simultaneous} shows that any such fixpoint is a sound
simultaneous omission.

\paragraph{Type arguments.}
The constructor $\mathit{cons} : \forall A{:}\mathsf{Type}.\ A \arr
\mathit{list}\,A \arr \mathit{list}\,A$ is translated as a ternary symbol
$\mathit{cons}(\alpha,x,l)$.  The type argument is recoverable from the
declared type of either term argument: $A$ occurs in the argument types $A$
and $\mathit{list}\,A$.  Erasing it --- replacing every occurrence
$\mathit{cons}(u,t,s)$ by $\mathit{cons}(t,s)$ uniformly in all premises,
axioms and the goal --- is sound: in the intended model the omitted argument
is a function of the types of the retained ones (\cref{lem:inference-unique}).
By contrast $\mathit{nil} : \forall A{:}\mathsf{Type}.\ \mathit{list}\,A$
has no term argument at all; its type argument is determined only by the
\emph{expected} type at each occurrence, which has no counterpart at a
first-order term position.  Erasing it conflates $\mathit{nil}\,\mathit{nat}$
with $\mathit{nil}\,\mathit{bool}$ as one constant, and
\cref{prop:tight-nil} turns this into an actual refutation of a goal that
fails in a model of the source theory --- the formal version of the classic
``$\mathit{card}\ \mathsf{UNIV}$'' unsoundness of
Meng and Paulson~\cite{MengPaulson2008}.  Finally, the two optimizations
interact: the cover $\mathit{Even}(x)$ forces $x : \mathit{nat}$ only
because the predicate symbol carries its full argument list.  Erasing type
arguments of \emph{predicates} destroys covers, and
\cref{prop:tight-pred-erasure} shows the combination is unsound; the
composition theorem (\cref{thm:composition}) therefore erases only function
and constructor symbols.

\paragraph{Proof architecture.}
All soundness statements are relative to the \emph{ground companion theory}
$\Th$ of the source environment $\Env$: the many-sorted first-order theory of
all ground type instances of the premises together with the structural
axioms of the inductive definitions.  \Cref{sec:fragment} defines $\Th$ and
shows that CIC proves every sentence we place in it, so that a model of
$\Th$ exists whenever the CIC environment is consistent.  Given any model
$\NN \models \Th$, \cref{sec:translation} constructs an untyped first-order
model $\MN$ by disjointly uniting (i) the ground types themselves, as
denotations of type expressions, (ii) tagged copies
$\{(\tau,a) \mid a \in \NN_\tau\}$ of the carriers, and (iii) an inductively
generated universe of \emph{junk} elements: formal application trees that
record a function symbol applied to arguments violating its type discipline.
Junk is what makes the model faithful to the untyped setting --- the
quantifiers of the translated problem really do range over ill-typed terms
--- and its \emph{free} (injective, non-confusable) construction is what
validates the unguarded injectivity and discrimination axioms that hammer
translations emit.  Soundness of each optimization then reduces to checking
that the transformed axioms remain true in $\MN$ (or in its erased variant
$\MNe$), and the final theorems all have the same shape: \emph{if the
optimized translation of $(\Env,\phi_0)$ is refutable, then $\phi_0$ holds
in every model of $\Th$.}
